\documentclass[11pt]{article}
\usepackage[margin=1in]{geometry}
\usepackage{amsmath,amssymb,amsthm}
\usepackage{booktabs}
\usepackage{hyperref}

\title{Moran Model for MERFISH Spatial Deconvolution:\\A Research Plan}
\author{Ho Lung Tsui}
\date{\today}

\begin{document}
\maketitle

%% ==================================================================
\section{The Spatial Deconvolution Task}
%% ==================================================================

MERFISH (Vizgen) provides single-cell resolution gene expression with spatial coordinates. However, many spatial transcriptomics platforms (e.g.\ Visium) only measure \emph{bulk} expression at each spatial spot---the sum across all cells at that location. \textbf{Spatial deconvolution} is the inverse problem: given the bulk measurement plus morphological side-information, recover the individual cell profiles.

\subsection{Precise Problem Statement}

\paragraph{Input (observed at test time).}
\begin{itemize}
  \item Bulk gene expression at a tissue spot: $X_0 \in \mathbb{Z}_{\geq 0}^{649}$, the sum of all cells' counts.
  \item DAPI nuclear images for each cell in the spot: $I_1, \ldots, I_N$, each $256 \times 256$ grayscale.
  \item The number of cells $N$ in the spot (known from nuclear segmentation).
\end{itemize}

\paragraph{Output (generated).}
\begin{itemize}
  \item Per-cell gene expression profiles: $\hat{x}_0^{(1)}, \ldots, \hat{x}_0^{(N)} \in \mathbb{Z}_{\geq 0}^{649}$.
\end{itemize}

\paragraph{Hard constraint.}
\[
  \sum_{i=1}^N \hat{x}_0^{(i)} = X_0.
\]
The generated cell profiles must sum exactly to the observed bulk measurement. This is enforced via IPF (iterative proportional fitting) rescaling followed by randomised rounding.

\paragraph{Why DAPI helps.} Nuclear morphology (size, shape, texture) correlates with cell type. The model uses each cell's DAPI image to infer its likely cell type, then assigns expression profiles accordingly---subject to the sum constraint.

\subsection{Evaluation (Against MERFISH Ground Truth)}

Since MERFISH provides the true single-cell profiles, we can evaluate directly:
\begin{enumerate}
  \item \textbf{Distributional metrics} (per spot): energy distance, sliced Wasserstein (100 projections), MMD with RBF kernel, MSE---comparing predicted vs.\ true cell profiles.
  \item \textbf{Cell type proportions}: Leiden clustering on the reference data at various resolutions (5, 10, 15, 20, 25, 30, 35 clusters). Labels are transferred to predicted cells via $k$-NN in PCA space. The cell type proportion (CTP) matrix per spot is compared.
  \item \textbf{Statistical moments}: mean, variance, skewness, kurtosis errors per gene; covariance Frobenius norm.
\end{enumerate}

Our Moran model must be evaluated with \textbf{exactly these metrics}, on \textbf{exactly this data}, to enable a direct comparison with Count Bridges.

%% ==================================================================
\section{How Count Bridges Solves This Task}
%% ==================================================================

\subsection{The Skellam Bridge: Per-Cell, Per-Gene, Independent}

Count Bridges treats each cell's each gene dimension as an \emph{independent} 1D birth-death bridge. For cell $i$, gene $j$:
\[
  x_t^{(i,j)} = x_0^{(i,j)} + B_t^{(i,j)} - D_t^{(i,j)},
\]
where $(B_t, D_t)$ are drawn from the bridge conditional $P(x_t \mid x_0, x_1)$, which has exact closed-form sampling via Binomial $+$ Hypergeometric $+$ Bessel slack. The source is i.i.d.\ noise: $x_1^{(i)} \sim |\mathcal{N}(0, 10)|$.

\textbf{There is no interaction between cells.} Cell $i$ and cell $j$ at the same spot are bridged completely independently. The only inter-cell coupling is the aggregation constraint, enforced post-hoc.

\subsection{Architecture: MultimodalUViT}

Each cell is processed independently by a 10-layer Vision Transformer (768-dim, 12 heads):
\begin{itemize}
  \item DAPI image $I_i$ ($256 \times 256$) $\to$ 256 patch embeddings via ViT patch embedding.
  \item Count vector $x_t^{(i)}$ (649-dim) $\to$ 4 patch embeddings via learned linear projection.
  \item Time token $+$ noise token $+$ 256 image patches $+$ 4 count patches $= 262$ tokens total.
  \item Output: predicted count vector $\hat{x}_0^{(i)} \in \mathbb{R}_{\geq 0}^{649}$ (Softplus ensures non-negativity).
\end{itemize}

The model sees one cell at a time. It cannot see other cells in the same spot.

\subsection{Training: E-M Loop}

Each epoch alternates two steps:
\begin{enumerate}
  \item \textbf{E-step} (inference): Sample $x_0$ from the current model via the reverse Skellam bridge, conditioned on $X_0$ via IPF rescaling. This produces ``pseudo-ground-truth'' cell profiles that are consistent with the bulk measurement.
  \item \textbf{M-step} (training): Train on $(x_0^{\text{inferred}}, x_1)$ pairs. The Skellam bridge produces intermediate states $(t, x_t, x_0)$. The energy score loss is computed at the \emph{group level} after aggregating cell predictions via the sparse matrix $A$:
  \[
    \mathcal{L} = \text{EnergyScore}\!\bigl(A \hat{x}_0,\; X_0\bigr),
  \]
  where $A \in \{0,1\}^{G \times C}$ maps cells to spots: $A_{g,c} = 1$ iff cell $c \in$ spot $g$.
\end{enumerate}

\subsection{Generation (Inference)}

At test time:
\begin{enumerate}
  \item Start from noise: $x_1^{(i)} \sim |\mathcal{N}(0, 10)|$ for each cell.
  \item Reverse through the Skellam bridge: at each step $k$, the model predicts $\hat{x}_0^{(i)}$ from $(x_t^{(i)}, t, I_i)$, then the bridge conditional gives $x_{t-1}^{(i)}$.
  \item At $t = 0$: apply IPF rescaling and randomised rounding to enforce $\sum_i \hat{x}_0^{(i)} = X_0$ exactly.
\end{enumerate}

Each cell is generated independently. The sum constraint is the only coordination.

\subsection{What This Means}

Two adjacent cells at the same spot---even if they are clearly the same cell type based on their nuclear morphology---receive \emph{completely independent} noise trajectories and model predictions. The only thing preventing wildly inconsistent profiles is the post-hoc IPF projection, which is a hard algebraic constraint, not a learned generative mechanism.

%% ==================================================================
\section{Count Bridges Benchmark Results on MERFISH}
%% ==================================================================

The evaluation compares three quantities per spot:
\begin{itemize}
  \item \textbf{Pred vs True}: Generated cell profiles vs.\ MERFISH ground truth. This is the primary metric---how accurate is the deconvolution?
  \item \textbf{Pred vs Spot Mean}: Generated profiles vs.\ the naive baseline of assigning every cell the spot mean $\bar{x}_0 = X_0 / N$. Measures whether the model captures \emph{any} cell-to-cell variability.
  \item \textbf{True vs Spot Mean}: True single-cell profiles vs.\ their own spot mean. This is the inherent biological variability---an oracle ceiling. No method can do better than the true cells against their own mean.
\end{itemize}

\subsection{Results: S1R1 (4{,}358 spots, 73{,}613 cells)}

\begin{center}
\begin{tabular}{lccc}
\toprule
Metric & Pred vs True & Pred vs Spot Mean & True vs Spot Mean \\
\midrule
Energy distance    & 8.891  & 42.903 & 41.717 \\
Sliced Wasserstein & 0.017  & 0.034  & 0.030 \\
MMD (RBF)          & 0.203  & 0.419  & 0.409 \\
\bottomrule
\end{tabular}
\end{center}

\subsection{Interpretation}

\begin{enumerate}
  \item \textbf{The model works well}: Pred vs True (8.9 energy distance) is much smaller than the inherent variability (41.7), meaning generated profiles are substantially closer to truth than to the naive spot-mean baseline.

  \item \textbf{Variability is captured}: Pred vs Spot Mean ($\approx 42.9$) is close to True vs Spot Mean ($\approx 41.7$), indicating Count Bridges produces roughly the correct \emph{amount} of cell-to-cell variability. The model does not under- or over-disperse relative to the data.

  \item \textbf{This is the target to beat}. Our Moran model must be evaluated with \emph{exactly} this three-way comparison on the same S1R1 data. If Moran's inter-cell coupling improves deconvolution, we should see Pred vs True decrease (closer to ground truth) while Pred vs Spot Mean remains close to True vs Spot Mean (preserving correct variability).
\end{enumerate}

\subsection{Additional Evaluation: Cell Type Proportions}

Beyond distributional metrics, Count Bridges evaluates cell type recovery:
\begin{enumerate}
  \item Leiden clustering on reference data at various resolutions (5, 10, 15, 20, 25, 30, 35 clusters).
  \item Labels transferred to predicted cells via $k$-NN ($k = 15$) in PCA space.
  \item Cell type proportion (CTP) matrix compared per spot: does the model assign the right mix of cell types?
\end{enumerate}

This is where the Moran model's DP prior (which naturally clusters cells into types) may provide an advantage over Count Bridges' independent generation.

%% ==================================================================
\section{What the Moran Model Changes}
%% ==================================================================

\subsection{Core Idea}

Instead of generating each cell independently, the Moran model generates all $N$ cells in a spot \emph{jointly}. Cells that are similar (spatially proximate, morphologically alike) influence each other through the resampling mechanism, producing correlated expression profiles as a mathematical consequence of the population dynamics.

\subsection{Framework}

The cells within each tissue spot form a Moran population:
\begin{itemize}
  \item \textbf{Individuals}: cells $\{1, \ldots, N\}$ within a spot.
  \item \textbf{Type of individual $i$}: gene expression profile $x^{(i)} \in \mathbb{Z}_{\geq 0}^{649}$.
  \item \textbf{Mutation}: per-gene immigration-death. Gene $j$ in cell $i$ has immigration rate $\beta$ and per-molecule death rate $\mu$, giving stationary distribution $\text{Poisson}(\beta/\mu)$ per gene.
  \item \textbf{Resampling}: cell $i$ copies cell $j$'s \emph{entire} 649-dim expression profile, at rate $\frac{\kappa}{N} K(i, j)$.
\end{itemize}

The generator of the Moran CTMC on the joint state $(x^{(1)}, \ldots, x^{(N)})$:
\begin{align}
  \mathcal{A}f &= \underbrace{\sum_{i=1}^N \mathcal{A}_{\text{mut}}^{(i)} f}_{\text{per-cell mutation}} + \underbrace{\frac{\kappa}{N} \sum_{i \neq j} K(i,j) \bigl[f(x^{i \leftarrow j}) - f(x)\bigr]}_{\text{Moran resampling}}.
\end{align}

\subsection{Choice of Kernel}

The resampling kernel $K(i,j)$ determines which cells influence each other. Three options, in order of preference:

\paragraph{Option A: Spatial kernel (recommended --- coordinates available!).}
$K(i, j) = \exp\!\bigl(-\|p_i - p_j\|^2 / 2h^2\bigr)$,
where $p_i = (\texttt{x\_um}_i, \texttt{y\_um}_i)$ are cell centroids in micrometres, available in the \texttt{S1R1.npz} files. Biologically motivated: spatial proximity implies shared lineage and microenvironment. This is the most natural choice and can be used immediately.

\paragraph{Option B: Uniform kernel within spot (simplest baseline).}
$K(i,j) = 1$ for all cells in the same spot. All cells are equally coupled. Useful for ablation: does the spatial structure of the kernel matter, or is any coupling sufficient?

\paragraph{Option C: UMAP/PCA kernel from reference expression.}
$K(i,j) = \exp\!\bigl(-\|u_i - u_j\|^2 / 2h^2\bigr)$,
where $u_i$ is the PCA or UMAP coordinate of cell $i$, computed from the \emph{reference} expression data. This captures transcriptomic similarity directly. The kernel is precomputed from ground truth and fixed during training, avoiding the chicken-and-egg problem of state-dependent rates.

\paragraph{Data availability.} The \texttt{S1R1.npz} files contain more fields than the dataset loader uses:
\begin{center}
\texttt{imgs, counts, spots, x\_um, y\_um, annotations, visium\_simulated, cell\_type\_counts, n\_cells}
\end{center}
\textbf{Spatial coordinates are available}: \texttt{x\_um} and \texttt{y\_um} give cell centroid positions in micrometres. The dataset loader (\texttt{merfish\_deconv.py}) only reads \texttt{imgs}, \texttt{counts}, and \texttt{spots}---it ignores the spatial coordinates entirely. This means a proper spatial kernel is immediately available without needing the raw Vizgen data.

\subsection{De Finetti Equilibrium as Prior}

With time scaling $s = 1/(1-t)$, the Moran process converges to its de Finetti equilibrium at $t = 1$:
\[
  \mu_N \;\xrightarrow{N \to \infty}\; \text{DP}(\theta, \nu),
\]
where $\theta = 2\gamma_{\text{mut}}/\kappa$ and $\nu$ is the stationary distribution of the per-cell mutation on $\mathbb{Z}_{\geq 0}^{649}$.

The DP prior \emph{naturally produces cell type clusters}: atoms of the DP are distinct expression profiles (cell types), shared by groups of cells, with $\theta$ controlling the expected number of distinct types. Compare this to Count Bridges' source $x_1^{(i)} \sim |\mathcal{N}(0, 10)|$ i.i.d., which has no type structure whatsoever.

\subsection{What Moran Provides That Count Bridges Does Not}

\begin{center}
\begin{tabular}{lcc}
\toprule
Property & Count Bridges & Moran \\
\midrule
Cell generation & Independent per cell & Joint (all cells in spot) \\
Inter-cell coupling & None (only IPF post-hoc) & Resampling during generation \\
Prior structure & i.i.d.\ noise & DP clusters (cell types) \\
Spatial coherence & Only via DAPI images & Explicit spatial/similarity kernel \\
Aggregation & Hard constraint only & Soft coupling $+$ hard constraint \\
\bottomrule
\end{tabular}
\end{center}

The central prediction: Moran should produce more \emph{internally consistent} cell populations per spot, where cells of the same type receive correlated expression profiles through the generative process itself, rather than relying solely on IPF projection.

%% ==================================================================
\section{Implementation Plan}
%% ==================================================================

\subsection{Phase 1: Two Moons $N$-Sweep (Proof of Concept)}

Validate the Moran particle framework on the discrete two moons dataset, using the same data and metrics as Count Bridges (value\_range $= 196$, MMD, Wasserstein, coverage, off-support).

The 50,000 data points are partitioned into populations of size $N$:
\begin{center}
\begin{tabular}{ccl}
\toprule
$N$ & Training populations & Note \\
\midrule
1 & 50,000 & Count Bridges limit ($\kappa = 0$) \\
10 & 5,000 & Light coupling \\
50 & 1,000 & Moderate coupling \\
100 & 500 & Strong coupling \\
\bottomrule
\end{tabular}
\end{center}

Sweep over $N$ and $\theta$. The model is a permutation-equivariant DeepSets denoiser with energy score loss. Implemented in \texttt{moran/experiment.py}.

\subsection{Phase 2: MERFISH Data Exploration}

The \texttt{S1R1.npz} file contains all needed fields. Key exploration tasks:
\begin{enumerate}
  \item Distribution of cells per spot (min, max, mean, histogram). From the notebooks: ${\sim}73{,}613$ cells across 4{,}358 spots in S1R1 (average ${\sim}17$).
  \item \textbf{Spatial structure within spots}: plot \texttt{x\_um}, \texttt{y\_um} for cells within individual spots. How clustered are they? What is the typical inter-cell distance? This determines the kernel bandwidth $h$.
  \item Gene expression structure: PCA spectrum, Leiden clustering at various resolutions (notebooks show 5--35 clusters). This calibrates $\theta$.
  \item Spot-to-spot variability: how heterogeneous are spots in cell type composition?
\end{enumerate}

\subsection{Phase 3: Moran Forward Process for 649-Dim Counts}

Implement the Moran CTMC for gene expression vectors:
\begin{itemize}
  \item \textbf{Mutation}: per-gene immigration-death with rates calibrated to match data statistics.
  \item \textbf{Resampling}: cell $i$ copies cell $j$'s entire 649-dim profile, weighted by kernel $K$.
  \item \textbf{Time scaling}: rates $\times\, 1/(1-t)^2$ to guarantee convergence to de Finetti.
  \item \textbf{Tau-leaping}: approximate simulation with ${\sim}30$ substeps per unit time.
\end{itemize}

\subsection{Phase 4: Architecture Adaptation}

Wrap the existing MultimodalUViT in a DeepSets framework:
\begin{enumerate}
  \item Per-cell UViT encoder (reuse existing architecture, including DAPI image processing).
  \item Mean-pool cell embeddings across the spot $\to$ global context vector.
  \item Per-cell decoder conditioned on global context.
\end{enumerate}

This adds inter-cell awareness with minimal modification to the existing architecture.

\subsection{Phase 5: Training and Evaluation}

\begin{itemize}
  \item Replace the Skellam bridge with the Moran forward/reverse in both E-step and M-step. Keep the E-M structure.
  \item Keep the aggregation constraint ($\sum_i \hat{x}_0^{(i)} = X_0$) and IPF rescaling.
  \item Evaluate with \textbf{exactly the same metrics} as Count Bridges: energy distance, Wasserstein, MMD, cell type proportions, moment errors.
  \item Additional Moran-specific analysis: do the DP clusters from the prior correspond to biologically meaningful cell types?
\end{itemize}

%% ==================================================================
\section{Key Open Questions}
%% ==================================================================

\begin{enumerate}
  \item \textbf{Does inter-cell coupling help?} The DAPI images already provide per-cell morphological information. The Moran coupling adds transcriptomic coherence among similar cells---is this redundant, or does it provide signal beyond what DAPI captures?

  \item \textbf{Optimal $\theta$?} Too small $\to$ all cells collapse to one type. Too large $\to$ no coupling (recovers Count Bridges). The biological expectation is ${\sim}5$--$20$ cell types per spot; $\theta$ should be calibrated accordingly.

  \item \textbf{Aggregation constraint interaction.} The DP prior naturally partitions cells into types. Does this make the IPF rescaling smoother (less disruptive) compared to Count Bridges' independent prior, where IPF must do all the coordination?

  \item \textbf{Computational cost.} The Moran forward is $O(N^2)$ per substep per spot (kernel computation). With $N \sim 18$ cells per spot, this is ${\sim}324$ kernel evaluations---cheap. The bottleneck is the 649-dimensional state space, not the population size.

  \item \textbf{Variable $N$ per spot.} Different spots have 3--50+ cells. The DeepSets architecture handles variable input sizes naturally. The sparse aggregation collation already manages this in Count Bridges.
\end{enumerate}

\end{document}
